% TOP_PROBLEM: {"id": 10000057, "title": "Limit shape in Poisson–Voronoi metrics over $\\ell_p$ planes", "category_id": 19, "category": {"id": 19, "name": "probability", "display_name": "Probability", "description": "Problems involving probability theory, stochastic processes, and random structures.", "slug": "probability", "order_index": 19, "created_at": "2026-07-31T15:26:25.671Z"}, "status": "open", "problem_number": "AMR-099-0057", "set_id": 13}
% FIRST_POSTED: 2026-10-01
% ENTRY_KIND: statement_only
% UnsolvedMath ID 10000057; source catalogue code AMR-099-0057
% !TeX program = lualatex
% Definition and sources only; this file is not an LLM solution attempt.
\documentclass[11pt,a4paper]{article}
\usepackage[margin=25mm]{geometry}
\usepackage{fontspec}
\setmainfont{lmroman10-regular.otf}[BoldFont=lmroman10-bold.otf,
 ItalicFont=lmroman10-italic.otf,BoldItalicFont=lmroman10-bolditalic.otf]
\usepackage{amsmath,amssymb,mathtools}
\usepackage{unicode-math}
\setmathfont{latinmodern-math.otf}
\AtBeginDocument{\let\setminus\smallsetminus}
\usepackage{xurl,hyperref}
\hypersetup{colorlinks=true,urlcolor=blue}
\setcounter{secnumdepth}{0}
\setlength{\parindent}{0pt}
\setlength{\parskip}{5pt}
\setlength{\emergencystretch}{3em}
\begin{document}
\section*{Limit shape in Poisson–Voronoi metrics over $\ell_p$ planes}\label{10000057}
{\small UnsolvedMath ID 10000057 · AMR-099-0057 · Probability · AMR Open Problem Lists}
\subsection{Definitions and mathematical statement}
Fix $p\in[1,\infty]$, and equip $\mathbb R^2$ with $\|x\|_p=(|x_1|^p+|x_2|^p)^{1/p}$ for finite $p$, or $\|x\|_\infty=\max(|x_1|,|x_2|)$. Let $\Pi$ be a homogeneous Poisson point process of intensity one with respect to planar Lebesgue area: counts in disjoint bounded sets are independent Poisson random variables with means equal to their areas.

For each nucleus $z\in\Pi$, form the closed Voronoi cell
\[
C_z=\{x:\|x-z\|_p\le\|x-w\|_p\text{ for all }w\in\Pi\}.
\]
Two cells are adjacent when they share a boundary portion of positive length. Their adjacency graph has all edges of length one. Choose the cell containing the origin, using any deterministic tie convention on the null event of a tie, and let $\mathcal B_n$ be the union of cells at graph distance at most $n$ from it.

Determine the deterministic large-scale shape $K_p$ of these graph-metric balls. In particular, establish the appropriate shape theorem and identify $K_p$ so that
\[
d_H(n^{-1}\mathcal B_n,K_p)\longrightarrow0
\quad\text{almost surely},
\]
where $d_H$ is Hausdorff distance in the Euclidean plane. A description via directional time constants should determine their dependence on direction, not merely rename the unknown limit.

The adjacency metric counts cell crossings rather than Euclidean or $\ell_p$ length. Consequently the $\ell_p$ metric used to construct cells does not by itself determine the final ball shape. Even its scale is fixed here by Poisson intensity one and unit graph edges; an equivalent answer may describe the shape up to an explicitly stated normalization.
\subsection{Short English statement}
What macroscopic shape is produced by counting crossings of Poisson–Voronoi cells built using an arbitrary planar lp metric?
\subsection{Sources}
\begin{itemize}
\item[S1] ULAM.AI and the UnsolvedMath Contributors, supplied problems.json, record 10000057 (AMR-099-0057), AMR Open Problem Lists. Catalogue identity and source transcription; CC BY 4.0.
\url{https://huggingface.co/datasets/ulamai/UnsolvedMath}.
\item[S2] Benjamini - Euclidean vs. Graph Metric, Question 3.2, PDF page 4. Original source indexed by AMR-099-0057.
\url{https://www.wisdom.weizmann.ac.il/~itai/erd100.pdf#page=4}.
\item[S3] I. Benjamini, Euclidean vs. Graph Metric, primary numbered question and conventions.
\url{https://arquivo.pt/noFrame/replay/20201231041538id_/http://www.wisdom.weizmann.ac.il/~itai/erd100.pdf}.
\end{itemize}
\end{document}
